\documentclass[../../main]{subfiles}
\begin{document}

\part{Coupled RSVP--Polyxan Dynamics and Semantic Galaxies}
\label{part:coupled-dynamics}

\chapter*{Introduction to Part III}
\addcontentsline{toc}{chapter}{Introduction to Part III}

Part~I established the continuous side of the semantic universe:  
the RSVP manifold $(\mathcal{M},g)$, its intrinsic geometry, the
scalar--vector--entropy fields $(\Phi,\mathbf{v},S)$, the action
principle, the renormalization structure, the mixed-type analytic
phases, and the stratified singularity formation that resolves into
elliptic attractors.  

Part~II constructed the discrete half:  
typed Polyxan hypergraphs, curvature and cohomology, spectral invariants,
rewrite systems, fibrations with connection, quine structures, and the
universal media quine $\mathfrak{U}$ with 
$H^k(\mathfrak{U})=0$ for all $k>0$.

\medskip

\noindent
\textbf{Part III brings these two worlds together.}  
It develops the coupled dynamics of RSVP fields and Polyxan hypergraphs,
showing that continuous semantic geometry and discrete algebraic
structure are not merely compatible but mutually generative.  
Every act of meaning is simultaneously:

\begin{enumerate}[label=\textbf{(\alph*)}]
\item a continuous field configuration $(\Phi,\mathbf{v},S)$ on the
      semantic manifold $\mathcal{M}$, and  
\item a discrete Polyxan hypergraph $\mathcal{G}$ whose atoms, tags,
      curvature, and cohomology encode the algebraic skeleton of the
      same meaning event.
\end{enumerate}

\noindent
These two representations are tied together by a single geometric object:

\begin{equation}
\label{eq:embedding}
\iota : \mathcal{G} \hookrightarrow \mathcal{M},
\end{equation}

a curvature- and entropy-controlled embedding that identifies each typed
hypergraph with a stratified region of the RSVP manifold.  
The embedding preserves type, curvature deficit, sheaf structure,
tag-cohomology, and entropy production.  
In particular, the continuous singularity set of the RSVP flow
(Chapter~\ref{ch:singularities}) corresponds exactly to the incidence
strata of $\mathcal{G}$ under $\iota$.

\medskip

\noindent
\textbf{The Master Lagrangian of Part III.}  
The core of this Part is the coupled action

\begin{equation}
\label{eq:master-lagrangian}
\mathcal{L}[\Phi,\mathbf{v},S;\mathcal{G},\iota]
  = \mathcal{L}_{\mathrm{RSVP}}[\Phi,\mathbf{v},S;g]
  + \mathcal{L}_{\mathrm{Polyx}}[\mathcal{G}]
  + \mathcal{L}_{\mathrm{int}}[\Phi,\mathbf{v},S;\mathcal{G},\iota].
\end{equation}

Here:

\begin{itemize}
\item $\mathcal{L}_{\mathrm{RSVP}}$ is the scalar--vector--entropy Lagrangian
      of Part~I;  
\item $\mathcal{L}_{\mathrm{Polyx}}$ is the spectral action of the
      normalized Laplacian, curvature deficit, and sheaf cohomology;  
\item $\mathcal{L}_{\mathrm{int}}$ encodes the geometric coupling:
      curvature–type matching, entropy exchange, quine-stability,
      and the back-reaction of discrete hypergraph curvature on the
      RSVP metric $g$.
\end{itemize}

Variation of~\eqref{eq:master-lagrangian} produces the coupled field
equations, governing the evolution of semantic matter in its continuous and
discrete forms.

\medskip

\noindent
\textbf{Semantic Galaxies.}  
Under joint RSVP--Polyxan flow, most initial conditions collapse into
stable composite structures:  
localized regions of high semantic density in $(\mathcal{M},g)$, coupled
to quine-stable hypergraphs with vanishing cohomology above degree~0.
These attractors are the \emph{semantic galaxies}.  
They are characterised by:

\begin{enumerate}[label=\textbf{(\roman*)}]
\item an RSVP vacuum region where
      $\nabla_a T^{ab}=0$,  
      $\mathrm{Ric}(g)=0$,  
      $\Delta S=0$;  
\item a Polyxan hypergraph $\mathcal{G}$ whose curvature deficit
      vanishes,  
      whose Laplacian has no spurious zero modes,  
      and whose cohomology is harmonic;  
\item an embedding $\iota$ whose image is a stratified minimal surface
      for the coupled energy functional.
\end{enumerate}

Semantic galaxies are thus fixed points of the joint RSVP--Polyxan
dynamics: the precise analogues of vacuum solutions in field theory and
stable quines in algebraic semantics.

\medskip

\noindent
\textbf{The Universal Ground State.}  
At the bottom of the coupled theory lies a unique global minimizer:

\begin{equation}
(\Phi,\mathbf{v},S,\mathcal{G},\iota) = 
(\Phi^\ast,0,S^\ast,\mathfrak{U},\iota^\ast).
\end{equation}

Here:
\begin{itemize}
\item $\mathfrak{U}$ is the universal media quine, carrying no positive-degree
      cohomology;  
\item $(\Phi^\ast,0,S^\ast)$ is the RSVP conformal vacuum, carrying minimal
      semantic energy;  
\item the embedding $\iota^\ast$ is flat, harmonic, and compatible with
      all curvature invariants.
\end{itemize}

This is the unique state where both discrete and continuous curvature vanish,
tag entropy is zero, and no semantic degrees of freedom propagate.  
It is the coupled theory’s analogue of the conformal fixed point of Part~I
and the cohomologically trivial quine of Part~II.

\medskip

\noindent
\textbf{The Unified Irreversibility Principle.}  
The joint system obeys a single Lyapunov structure:

\[
c_{\mathrm{total}}(t)
  := c_{\mathrm{RSVP}}(t) + c_{\mathrm{Polyx}}(t)
  \quad\text{is strictly decreasing in time},
\]

with equality \emph{only} at a semantic galaxy.
This is the unification of:
\begin{itemize}
\item the RSVP $c$-theorem of Chapter~\ref{ch:rsvp-renormalization}, and  
\item the Polyxan $c$-theorem of Chapter~\ref{ch:spectral-semantics}.
\end{itemize}

Irreversibility is therefore not merely dynamical but structural:
semantic curvature and semantic entropy can only decrease under the
allowed transformations.  
All information not stabilised into a semantic galaxy is smoothed away.

\medskip

\noindent
\textbf{Structure of Part III.}  
The chapters that follow carry out the programme sketched above:
the construction of the master embedding, the coupled Euler--Lagrange
equations, the existence and uniqueness of RSVP--Polyxan flow, the
classification of semantic galaxies, and the proof that $\mathfrak{U}$
is the universal ground state.

Part III is the bridge between geometric physics and algebraic semantics.
It is here that the hyperstructure becomes a single, unified universe.

\end{document}

